% ===============================================================
% Chapter 1 — Introduction
% ===============================================================
\chapter{Introduction}\label{cap:intro}

% --- Section 1.1 ---
\section{The Case for Adaptive IoT Security}\label{sec:case_for_adaptive_security}
The Internet of Things (IoT) connects billions of embedded devices --- smart cameras, industrial sensors, medical monitors, and home appliances --- into a single, globally addressable network. The number of IoT devices worldwide is forecast to more than double from 19.8 billion in 2025 to over 40.6 billion by 2034 \cite{Statista2025}, an exponential growth that has produced a vast and multifaceted attack surface. The financial ramifications are staggering: global damages from cybercrime are projected to reach USD 12.2 trillion annually by 2031 \cite{Morgan2025}. Malicious actors increasingly exploit insecure devices to steal data, disrupt services, or form large-scale botnets for Distributed Denial of Service (DDoS) attacks \cite{Neto2023,Tawalbeh2020}.

Securing this landscape presents unique and formidable challenges that render traditional IT security paradigms insufficient. First, most IoT devices are severely \textbf{resource-constrained}, possessing minimal computational power, memory, and energy, which precludes the use of intensive security software like traditional firewalls or antivirus agents \cite{Chen2021}. Second, IoT ecosystems are defined by extreme \textbf{heterogeneity}, comprising countless devices using disparate protocols and operating systems, which makes uniform security policy enforcement nearly impossible \cite{Mavroeidis2023}. This complexity is amplified by the large-scale, dynamic nature of these networks and the physical vulnerability of devices often deployed in unsecured locations \cite{AlGaradi2020}.

Consequently, conventional security tools are fundamentally ill-equipped for this environment. \textbf{Signature-based Intrusion Detection Systems (IDS)}, the cornerstone of many legacy systems, rely on static databases of known attack patterns. This approach is purely reactive and fails against novel or \textbf{zero-day attacks} --- threats for which no signature yet exists \cite{Khraisat2019,Yang2024}. While \textbf{anomaly-based IDS} attempt to overcome this by baselining ``normal'' behaviour, the dynamic and heterogeneous nature of IoT traffic makes establishing a reliable baseline exceptionally difficult, leading to high \textbf{false positive rates} that can disrupt critical functions \cite{Ibitoye2021,Tharewal2022}.

The core limitation of these paradigms is their lack of adaptability. They are built on static models and cannot autonomously evolve. This creates a critical need for a new security paradigm: one that is forward-looking, data-driven, and capable of autonomous decision-making. \textbf{Deep Reinforcement Learning (DRL)}, a field of machine learning focused on goal-oriented learning from interaction, has emerged as a powerful framework for such a system \cite{Chen2021}. A DRL agent learns an optimal decision-making strategy, or \textbf{policy}, through trial-and-error interaction with its environment. Unlike static systems, a DRL agent can generalise from its experience to respond to novel threats and can adapt its policy as new adversarial tactics emerge, making it well-suited for the non-stationary nature of IoT security \cite{Sutton2018,Uprety2021}. By framing cyber-defense as a sequential decision-making problem, DRL provides a robust mathematical foundation for creating autonomous security systems that can learn to outmaneuver attackers \cite{Gueriani2023}.

% --- Section 1.2 ---
\section{Problem Statement and Contributions}\label{sec:proposal_and_contributions}
The preceding analysis culminates in a clear research problem: existing IoT security frameworks are not simultaneously \textbf{adaptive}, \textbf{data-driven}, and \textbf{rigorously benchmarked} against deployable reference policies under a unified evaluation contract. Reactive signature-based systems cannot generalise to novel attacks, supervised classifiers do not act, and prior DRL-for-IoT work seldom reports head-to-head comparisons against the full set of deployable baselines a security operator has at hand --- random, always-block, and a supervised stage classifier composed with a recommended-action mapping.

This dissertation directly addresses that gap by designing, implementing, and rigorously evaluating a closed-loop adaptive defense framework for IoT networks that combines LSTM-based kill-chain stage modeling with Deep Reinforcement Learning, validated end-to-end on the contemporary CICIoT2023 dataset. Concretely, this work makes the following contributions:

\begin{enumerate}
    \item \textbf{A kill-chain-aware adaptive defense framework.} A two-stage architecture combining (i) a Long Short-Term Memory (LSTM) red-team episode generator trained on stage-labelled CICIoT2023 sequences, which drives a \texttt{RealisationEngine} that samples real feature rows per kill-chain stage, with (ii) a custom Gymnasium environment exposing a five-action defender (observe, log, throttle, block, isolate) under a stage-action proportionality reward inspired by the IoTWarden recommended-action mapping~\cite{Alam2024}. The attacker model is bounded by a \textbf{monotonic-attacker assumption}: the LSTM generates upper-triangular stage transitions in which the kill chain advances but never retreats. The framework is developed and validated on the \textbf{CICIoT2023 dataset}~\cite{Neto2023}, projected onto a five-stage Cyber Kill Chain (benign, recon, access, maneuver, impact) with a documented per-class mapping, with out-of-distribution classes held out disjointly from training.

    \item \textbf{A reward-mis-specification analysis for kill-chain-aware RL.} A naive contract that terminates the episode at the first \textsc{impact} step ({\small\texttt{impact\_is\_terminal=True}}) causes the agent to learn a cocked-trigger strategy that scores deceptively high on validation but fails to achieve genuine threat mitigation at benchmark scale. Identifying and correcting this mis-specification --- by allowing the agent one additional decision turn after \textsc{impact} ({\small\texttt{impact\_is\_terminal=False}}) --- is a primary contribution of this work. The corrected primary contract is the basis for all reported benchmark numbers.

    \item \textbf{A comparative DRL benchmark with a ${\sim}146{\times}$ latency advantage over the supervised deployable baseline.} DQN, PPO, and A2C are trained over ten seeds under the primary contract and evaluated on a held-out balanced test split (300 episodes per policy) against four reference policies: random, always-observe, always-block, and a supervised RandomForest stage classifier composed with the recommended-action mapping (RF-Acting). A rule-based policy with free access to the true attack stage is reported as an \emph{oracle ceiling} (\OracleCeiling{}). The best deployable trained agent (A2C, \BestAgentReward{}) captures \textbf{\OracleCapturePct{}\,\%} of that ceiling with non-overlapping bootstrap confidence intervals against every non-RL baseline. Among deployable policies the comparison is a latency--reward trade-off: trained RL agents achieve inference at ${\sim}0.095$\,ms per decision versus RF-Acting at 13.83\,ms --- a ${\sim}146{\times}$ advantage --- at a reward cost of approximately 10\,\%. A caveat applies: the best-reward trained agent (A2C) also carries the highest benign false-positive rate (11.5\,\%), while DQN offers a lower-FPR operating point (6.1\,\%) at marginally lower reward.

    \item \textbf{Reward-component and out-of-distribution ablations.} A 12-cell sparse one-at-a-time sweep over five reward coefficients $\times$ $\{0.5\times, 1\times, 2\times\}$ plus the binary termination-contract flag characterises the reward landscape. No single linear perturbation closes the residual gap; the structural \texttt{impact\_is\_terminal} change is the highest-impact lever, which in an ablation probe (PPO-only, $n=30$ episodes) raises the mitigated-impact rate from 0.153 to 0.840. At full benchmark scale ($n=300$, 10 seeds) the primary-contract agents achieve mitigated-impact rates of 0.26--0.32, a genuine improvement over the mis-specified baseline (0.153). A separate evaluation against four held-out out-of-distribution attack classes yields the pre-registered finding that the trained policy is \emph{robust to}, but not better at, the class on which the supervised stage detector is structurally blind (recall\,$=\,0.001$).

    \item \textbf{An audit-first reproducibility protocol.} Every figure ships a \texttt{manifest.json} that pins the SHA-256 hashes of every input artefact (data splits, trained models, scaler, episode roll-outs) and the producing Git commit. All results are reproducible via a SHA-256-verified artifact chain. The protocol is documented in Appendix~\ref{app:repro}.
\end{enumerate}

The remainder of this dissertation is structured as follows. Chapter~\ref{cap:background} provides a detailed review of background concepts in IoT security, the Cyber Kill Chain, and reinforcement learning, and situates this work against related research. Chapter~\ref{cap:methodology} details the proposed framework architecture, including the kill-chain projection of CICIoT2023, the LSTM red team and stage detector, the Markov Decision Process underlying the environment, the reward function, the training protocol, and the baseline policies. Chapter~\ref{cap:results} presents the empirical results across all phases of the pipeline and discusses the headline findings, including the oracle-ceiling reframe, the reward-component ablation, and the out-of-distribution robustness story. Chapter~\ref{cap:conclusion} summarises the contributions, discusses limitations and threats to validity, and outlines future work.
